% ============================================================================
% Cover Letter for submission to Computers in Biology and Medicine (CBM)
% Paper: When Does Temperature Scaling Pay Off in Cross-Corpus ECG Transfer?
%        A Multi-Seed Calibration Boundary Study
% File: paper/cover_letter.tex
% ============================================================================

\documentclass[11pt,letterpaper]{letter}

\usepackage[utf8]{inputenc}
\usepackage[margin=1in]{geometry}
\usepackage{amsmath,amssymb}

\signature{Zero Token Lab\\Deployment Science Group\\Shouguang, China}
\address{Zero Token Lab\\Deployment Science Group\\Shouguang, China}

\begin{document}

\begin{letter}{Editor-in-Chief\\
\emph{Computers in Biology and Medicine}\\
Editorial Office}

\opening{Dear Editor,}

We are pleased to submit our original research article entitled
``\textbf{When Does Temperature Scaling Pay Off in Cross-Corpus ECG
Transfer? A Multi-Seed Calibration Boundary Study}'' for your
consideration as an \textbf{Original Research Article} in
\emph{Computers in Biology and Medicine}. This manuscript has not been
published elsewhere and is not under consideration by any other journal.

\bigskip

\textbf{Summary of contributions.}
To our knowledge, this is the first study to systematically
characterize both the payoff and the boundary conditions of
Temperature Scaling (TS) as a post-hoc recalibration method under
cross-corpus ECG transfer. We pre-registered and executed a multi-seed
calibration boundary study comprising \textbf{60 experiments}
(6 directed transfer pairs $\times$ 2 architectures
$\times$ 5 seeds, fully balanced) on PTB-XL, Chapman-Shaoxing, and CPSC2018+2019,
with patient-level cluster paired bootstrap ($B=10{,}000$; percentile
operating CI with BCa sensitivity on 3 representative experiments) and
eight recalibration methods. TS yields a positive out-of-distribution
(OOD) calibration benefit
$\Delta\text{ECE}_{\text{OOD}}=\text{ECE}_{raw}^{OOD}
-\text{ECE}_{TS}^{OOD}$ in \textbf{51/60 = 85.0\%} of seed
experiments (InceptionTime 27/30, ResNet1D 24/30), while the
remaining \textbf{9 counter-examples are transparently disclosed}
and mechanistically explained case by case (six of nine co-occur with
low OOD accuracy; seven of nine with a large ID--OOD accuracy gap).
The paper delivers
\textbf{four independent contributions}:
\textbf{C1} (OOD benefit quantification, 51/60 seed support),
\textbf{C2} (ID boundary condition with the pre-registered branch
triggered: only 29/60 = 48.3\% of cells have CIs containing zero,
below the $\geq$80\% criterion, so the ``ID nonzero boundary'' branch
applies and C1 is reported with this qualifier),
\textbf{C3} (27-cell Shapley decomposition validation,
$n=20{,}000$), and
\textbf{C4} (method-selection boundary: TS robust vs.\ EM prior and
Matrix scaling weakest, governed by a prior-likelihood mismatch
mechanism).

\bigskip

\textbf{Why this journal.}
Cross-corpus ECG transfer calibration is a critical bottleneck for
clinical deployment of deep ECG models across institutions with
different acquisition protocols, lead sets, and label vocabularies.
Our methodological contribution---an explicit boundary-condition
characterization of when TS pays off, combined with transparent
reporting of nine counter-examples and a pre-registered failure
branch---aligns directly with the journal's scope on biomedical
signal processing, clinical machine learning reliability, and
reproducible computational methods in medicine.

\bigskip

\textbf{Transparency statement.}
All hypotheses in this manuscript are tagged
\emph{exploratory + robustness validation}, not confirmatory,
consistent with the open-science best practices of
Nosek et al.\ (2018). The nine counter-examples (15.0\% of the
experiment grid) are disclosed in full in the main
text and supplements; no non-supporting result is omitted or
post-hoc excluded. The pre-registration protocol v2.1-A1, its hash
manifest, and revision A1 are included with the submission
as supplementary files;
public OSF archival is in progress and the hash manifest allows
independent verification of the snapshot timing. All deviations from
the original protocol are recorded in the protocol revision log
(protocol \S11.5, amendment row A2).
The source code, experiment artifacts, and data preparation
scripts are provided with the submission as supplementary material;
a public repository will accompany acceptance.
A five-item honesty declaration (adversarial audit product) is
included in the manuscript and archived with the protocol.

\bigskip

\textbf{Note on pre-registration design.}
This study adopts an \emph{exploratory + robustness validation}
design rather than a strictly confirmatory design. The
pre-registration protocol was made public before the multi-seed
validation was unblinded. The formal revision A1 records a
transparent re-location of the primary endpoint from the original
``ID$\rightarrow$OOD decay'' to ``OOD benefit + boundary
conditions''. This re-location was triggered not by the main-test
results but by an independent exploratory pilot
($\Delta\text{ECE}_{\text{ID}}\approx 0$), and it is consistent with
the empirically observed ID--OOD asymmetry (Guo et al.\ 2017; Ovadia
et al.\ 2019); we emphasize this asymmetry is empirical, not a
theorem. Following Lakens (2021) and Nosek et al.\ (2018), this
constitutes a
\textbf{transparent re-location, not HARKing}: the pilot results
do not enter the main test report, and the revision direction is
consistent with the registered empirical asymmetry rather than
result-driven. We disclose this design choice at the outset so that
reviewers and editors can evaluate the evidence under the correct
inferential framework.

\bigskip

We believe this work makes a significant methodological
contribution to the field of cross-corpus ECG calibration,
providing an actionable answer to ``when does TS pay off'' with
multi-seed robustness evidence, transparent counter-example
reporting, and a pre-registered boundary-condition analysis. We
look forward to your feedback.

\closing{Sincerely,}

\end{letter}

\end{document}